\documentclass[11pt]{article}

\usepackage[breakable]{tcolorbox}
\usepackage{parskip} % Stop auto-indenting (to mimic markdown behaviour)
\usepackage{ctex}
\usepackage[table]{xcolor}

% Basic figure setup, for now with no caption control since it's done
% automatically by Pandoc (which extracts ![](path) syntax from Markdown).
\usepackage{graphicx}
% Keep aspect ratio if custom image width or height is specified
\setkeys{Gin}{keepaspectratio}
% Maintain compatibility with old templates. Remove in nbconvert 6.0
\let\Oldincludegraphics\includegraphics
% Ensure that by default, figures have no caption (until we provide a
% proper Figure object with a Caption API and a way to capture that
% in the conversion process - todo).
\usepackage{caption}
\DeclareCaptionFormat{nocaption}{}
\captionsetup{format=nocaption,aboveskip=0pt,belowskip=0pt}

\usepackage{float}
\floatplacement{figure}{H} % forces figures to be placed at the correct location
\usepackage{xcolor} % Allow colors to be defined
\usepackage{enumerate} % Needed for markdown enumerations to work
\usepackage{geometry} % Used to adjust the document margins
\usepackage{amsmath} % Equations
\usepackage{amssymb} % Equations
\usepackage{textcomp} % defines textquotesingle
% Hack from http://tex.stackexchange.com/a/47451/13684:
\AtBeginDocument{%
	\def\PYZsq{\textquotesingle}% Upright quotes in Pygmentized code
}
\usepackage{upquote} % Upright quotes for verbatim code
\usepackage{eurosym} % defines \euro

\usepackage{iftex}
\ifPDFTeX
\usepackage[T1]{fontenc}
\IfFileExists{alphabeta.sty}{
	\usepackage{alphabeta}
}{
	\usepackage[mathletters]{ucs}
	\usepackage[utf8x]{inputenc}
}
\else
\usepackage{fontspec}
\usepackage{unicode-math}
\fi

\usepackage{fancyvrb} % verbatim replacement that allows latex
\usepackage{grffile} % extends the file name processing of package graphics
% to support a larger range
\makeatletter % fix for old versions of grffile with XeLaTeX
\@ifpackagelater{grffile}{2019/11/01}
{
	% Do nothing on new versions
}
{
	\def\Gread@@xetex#1{%
		\IfFileExists{"\Gin@base".bb}%
		{\Gread@eps{\Gin@base.bb}}%
		{\Gread@@xetex@aux#1}%
	}
}
\makeatother
\usepackage[Export]{adjustbox} % Used to constrain images to a maximum size
\adjustboxset{max size={0.9\linewidth}{0.9\paperheight}}

% The hyperref package gives us a pdf with properly built
% internal navigation ('pdf bookmarks' for the table of contents,
% internal cross-reference links, web links for URLs, etc.)
\usepackage{hyperref}
% The default LaTeX title has an obnoxious amount of whitespace. By default,
% titling removes some of it. It also provides customization options.
\usepackage{titling}
\usepackage{longtable} % longtable support required by pandoc >1.10
\usepackage{booktabs}  % table support for pandoc > 1.12.2
\usepackage{array}     % table support for pandoc >= 2.11.3
\usepackage{calc}      % table minipage width calculation for pandoc >= 2.11.1
\usepackage[inline]{enumitem} % IRkernel/repr support (it uses the enumerate* environment)
\usepackage[normalem]{ulem} % ulem is needed to support strikethroughs (\sout)
% normalem makes italics be italics, not underlines
\usepackage{soul}      % strikethrough (\st) support for pandoc >= 3.0.0
\usepackage{mathrsfs}


\usepackage{geometry}
\usepackage{float}
\usepackage{multicol}
\usepackage{subfigure}
\usepackage{makecell}

% Colors for the hyperref package
\definecolor{urlcolor}{rgb}{0,.145,.698}
\definecolor{linkcolor}{rgb}{.71,0.21,0.01}
\definecolor{citecolor}{rgb}{.12,.54,.11}

% ANSI colors
\definecolor{ansi-black}{HTML}{3E424D}
\definecolor{ansi-black-intense}{HTML}{282C36}
\definecolor{ansi-red}{HTML}{E75C58}
\definecolor{ansi-red-intense}{HTML}{B22B31}
\definecolor{ansi-green}{HTML}{00A250}
\definecolor{ansi-green-intense}{HTML}{007427}
\definecolor{ansi-yellow}{HTML}{DDB62B}
\definecolor{ansi-yellow-intense}{HTML}{B27D12}
\definecolor{ansi-blue}{HTML}{208FFB}
\definecolor{ansi-blue-intense}{HTML}{0065CA}
\definecolor{ansi-magenta}{HTML}{D160C4}
\definecolor{ansi-magenta-intense}{HTML}{A03196}
\definecolor{ansi-cyan}{HTML}{60C6C8}
\definecolor{ansi-cyan-intense}{HTML}{258F8F}
\definecolor{ansi-white}{HTML}{C5C1B4}
\definecolor{ansi-white-intense}{HTML}{A1A6B2}
\definecolor{ansi-default-inverse-fg}{HTML}{FFFFFF}
\definecolor{ansi-default-inverse-bg}{HTML}{000000}

% common color for the border for error outputs.
\definecolor{outerrorbackground}{HTML}{FFDFDF}

% commands and environments needed by pandoc snippets
% extracted from the output of `pandoc -s`
\providecommand{\tightlist}{%
	\setlength{\itemsep}{0pt}\setlength{\parskip}{0pt}}
\DefineVerbatimEnvironment{Highlighting}{Verbatim}{commandchars=\\\{\}}
% Add ',fontsize=\small' for more characters per line
\newenvironment{Shaded}{}{}
\newcommand{\KeywordTok}[1]{\textcolor[rgb]{0.00,0.44,0.13}{\textbf{{#1}}}}
\newcommand{\DataTypeTok}[1]{\textcolor[rgb]{0.56,0.13,0.00}{{#1}}}
\newcommand{\DecValTok}[1]{\textcolor[rgb]{0.25,0.63,0.44}{{#1}}}
\newcommand{\BaseNTok}[1]{\textcolor[rgb]{0.25,0.63,0.44}{{#1}}}
\newcommand{\FloatTok}[1]{\textcolor[rgb]{0.25,0.63,0.44}{{#1}}}
\newcommand{\CharTok}[1]{\textcolor[rgb]{0.25,0.44,0.63}{{#1}}}
\newcommand{\StringTok}[1]{\textcolor[rgb]{0.25,0.44,0.63}{{#1}}}
\newcommand{\CommentTok}[1]{\textcolor[rgb]{0.38,0.63,0.69}{\textit{{#1}}}}
\newcommand{\OtherTok}[1]{\textcolor[rgb]{0.00,0.44,0.13}{{#1}}}
\newcommand{\AlertTok}[1]{\textcolor[rgb]{1.00,0.00,0.00}{\textbf{{#1}}}}
\newcommand{\FunctionTok}[1]{\textcolor[rgb]{0.02,0.16,0.49}{{#1}}}
\newcommand{\RegionMarkerTok}[1]{{#1}}
\newcommand{\ErrorTok}[1]{\textcolor[rgb]{1.00,0.00,0.00}{\textbf{{#1}}}}
\newcommand{\NormalTok}[1]{{#1}}

% Additional commands for more recent versions of Pandoc
\newcommand{\ConstantTok}[1]{\textcolor[rgb]{0.53,0.00,0.00}{{#1}}}
\newcommand{\SpecialCharTok}[1]{\textcolor[rgb]{0.25,0.44,0.63}{{#1}}}
\newcommand{\VerbatimStringTok}[1]{\textcolor[rgb]{0.25,0.44,0.63}{{#1}}}
\newcommand{\SpecialStringTok}[1]{\textcolor[rgb]{0.73,0.40,0.53}{{#1}}}
\newcommand{\ImportTok}[1]{{#1}}
\newcommand{\DocumentationTok}[1]{\textcolor[rgb]{0.73,0.13,0.13}{\textit{{#1}}}}
\newcommand{\AnnotationTok}[1]{\textcolor[rgb]{0.38,0.63,0.69}{\textbf{\textit{{#1}}}}}
\newcommand{\CommentVarTok}[1]{\textcolor[rgb]{0.38,0.63,0.69}{\textbf{\textit{{#1}}}}}
\newcommand{\VariableTok}[1]{\textcolor[rgb]{0.10,0.09,0.49}{{#1}}}
\newcommand{\ControlFlowTok}[1]{\textcolor[rgb]{0.00,0.44,0.13}{\textbf{{#1}}}}
\newcommand{\OperatorTok}[1]{\textcolor[rgb]{0.40,0.40,0.40}{{#1}}}
\newcommand{\BuiltInTok}[1]{{#1}}
\newcommand{\ExtensionTok}[1]{{#1}}
\newcommand{\PreprocessorTok}[1]{\textcolor[rgb]{0.74,0.48,0.00}{{#1}}}
\newcommand{\AttributeTok}[1]{\textcolor[rgb]{0.49,0.56,0.16}{{#1}}}
\newcommand{\InformationTok}[1]{\textcolor[rgb]{0.38,0.63,0.69}{\textbf{\textit{{#1}}}}}
\newcommand{\WarningTok}[1]{\textcolor[rgb]{0.38,0.63,0.69}{\textbf{\textit{{#1}}}}}
\makeatletter
\newsavebox\pandoc@box
\newcommand*\pandocbounded[1]{%
	\sbox\pandoc@box{#1}%
	% scaling factors for width and height
	\Gscale@div\@tempa\textheight{\dimexpr\ht\pandoc@box+\dp\pandoc@box\relax}%
	\Gscale@div\@tempb\linewidth{\wd\pandoc@box}%
	% select the smaller of both
	\ifdim\@tempb\p@<\@tempa\p@
	\let\@tempa\@tempb
	\fi
	% scaling accordingly (\@tempa < 1)
	\ifdim\@tempa\p@<\p@
	\scalebox{\@tempa}{\usebox\pandoc@box}%
	% scaling not needed, use as it is
	\else
	\usebox{\pandoc@box}%
	\fi
}
\makeatother

% Define a nice break command that doesn't care if a line doesn't already
% exist.
\def\br{\hspace*{\fill} \\* }
% Math Jax compatibility definitions
\def\gt{>}
\def\lt{<}
\let\Oldtex\TeX
\let\Oldlatex\LaTeX
\renewcommand{\TeX}{\textrm{\Oldtex}}
\renewcommand{\LaTeX}{\textrm{\Oldlatex}}
% Document parameters
% Document title
\title{\textbf{{\huge 机器视觉-实践项目4}}}
\vspace{2em}\author{国嘉通\thanks{Tel.+86-166 8130 6085,\ E-mail:guojiatom2006@outlook.com}} 

% Pygments definitions
\makeatletter
\def\PY@reset{\let\PY@it=\relax \let\PY@bf=\relax%
	\let\PY@ul=\relax \let\PY@tc=\relax%
	\let\PY@bc=\relax \let\PY@ff=\relax}
\def\PY@tok#1{\csname PY@tok@#1\endcsname}
\def\PY@toks#1+{\ifx\relax#1\empty\else%
	\PY@tok{#1}\expandafter\PY@toks\fi}
\def\PY@do#1{\PY@bc{\PY@tc{\PY@ul{%
				\PY@it{\PY@bf{\PY@ff{#1}}}}}}}
\def\PY#1#2{\PY@reset\PY@toks#1+\relax+\PY@do{#2}}

\@namedef{PY@tok@w}{\def\PY@tc##1{\textcolor[rgb]{0.73,0.73,0.73}{##1}}}
\@namedef{PY@tok@c}{\let\PY@it=\textit\def\PY@tc##1{\textcolor[rgb]{0.24,0.48,0.48}{##1}}}
\@namedef{PY@tok@cp}{\def\PY@tc##1{\textcolor[rgb]{0.61,0.40,0.00}{##1}}}
\@namedef{PY@tok@k}{\let\PY@bf=\textbf\def\PY@tc##1{\textcolor[rgb]{0.00,0.50,0.00}{##1}}}
\@namedef{PY@tok@kp}{\def\PY@tc##1{\textcolor[rgb]{0.00,0.50,0.00}{##1}}}
\@namedef{PY@tok@kt}{\def\PY@tc##1{\textcolor[rgb]{0.69,0.00,0.25}{##1}}}
\@namedef{PY@tok@o}{\def\PY@tc##1{\textcolor[rgb]{0.40,0.40,0.40}{##1}}}
\@namedef{PY@tok@ow}{\let\PY@bf=\textbf\def\PY@tc##1{\textcolor[rgb]{0.67,0.13,1.00}{##1}}}
\@namedef{PY@tok@nb}{\def\PY@tc##1{\textcolor[rgb]{0.00,0.50,0.00}{##1}}}
\@namedef{PY@tok@nf}{\def\PY@tc##1{\textcolor[rgb]{0.00,0.00,1.00}{##1}}}
\@namedef{PY@tok@nc}{\let\PY@bf=\textbf\def\PY@tc##1{\textcolor[rgb]{0.00,0.00,1.00}{##1}}}
\@namedef{PY@tok@nn}{\let\PY@bf=\textbf\def\PY@tc##1{\textcolor[rgb]{0.00,0.00,1.00}{##1}}}
\@namedef{PY@tok@ne}{\let\PY@bf=\textbf\def\PY@tc##1{\textcolor[rgb]{0.80,0.25,0.22}{##1}}}
\@namedef{PY@tok@nv}{\def\PY@tc##1{\textcolor[rgb]{0.10,0.09,0.49}{##1}}}
\@namedef{PY@tok@no}{\def\PY@tc##1{\textcolor[rgb]{0.53,0.00,0.00}{##1}}}
\@namedef{PY@tok@nl}{\def\PY@tc##1{\textcolor[rgb]{0.46,0.46,0.00}{##1}}}
\@namedef{PY@tok@ni}{\let\PY@bf=\textbf\def\PY@tc##1{\textcolor[rgb]{0.44,0.44,0.44}{##1}}}
\@namedef{PY@tok@na}{\def\PY@tc##1{\textcolor[rgb]{0.41,0.47,0.13}{##1}}}
\@namedef{PY@tok@nt}{\let\PY@bf=\textbf\def\PY@tc##1{\textcolor[rgb]{0.00,0.50,0.00}{##1}}}
\@namedef{PY@tok@nd}{\def\PY@tc##1{\textcolor[rgb]{0.67,0.13,1.00}{##1}}}
\@namedef{PY@tok@s}{\def\PY@tc##1{\textcolor[rgb]{0.73,0.13,0.13}{##1}}}
\@namedef{PY@tok@sd}{\let\PY@it=\textit\def\PY@tc##1{\textcolor[rgb]{0.73,0.13,0.13}{##1}}}
\@namedef{PY@tok@si}{\let\PY@bf=\textbf\def\PY@tc##1{\textcolor[rgb]{0.64,0.35,0.47}{##1}}}
\@namedef{PY@tok@se}{\let\PY@bf=\textbf\def\PY@tc##1{\textcolor[rgb]{0.67,0.36,0.12}{##1}}}
\@namedef{PY@tok@sr}{\def\PY@tc##1{\textcolor[rgb]{0.64,0.35,0.47}{##1}}}
\@namedef{PY@tok@ss}{\def\PY@tc##1{\textcolor[rgb]{0.10,0.09,0.49}{##1}}}
\@namedef{PY@tok@sx}{\def\PY@tc##1{\textcolor[rgb]{0.00,0.50,0.00}{##1}}}
\@namedef{PY@tok@m}{\def\PY@tc##1{\textcolor[rgb]{0.40,0.40,0.40}{##1}}}
\@namedef{PY@tok@gh}{\let\PY@bf=\textbf\def\PY@tc##1{\textcolor[rgb]{0.00,0.00,0.50}{##1}}}
\@namedef{PY@tok@gu}{\let\PY@bf=\textbf\def\PY@tc##1{\textcolor[rgb]{0.50,0.00,0.50}{##1}}}
\@namedef{PY@tok@gd}{\def\PY@tc##1{\textcolor[rgb]{0.63,0.00,0.00}{##1}}}
\@namedef{PY@tok@gi}{\def\PY@tc##1{\textcolor[rgb]{0.00,0.52,0.00}{##1}}}
\@namedef{PY@tok@gr}{\def\PY@tc##1{\textcolor[rgb]{0.89,0.00,0.00}{##1}}}
\@namedef{PY@tok@ge}{\let\PY@it=\textit}
\@namedef{PY@tok@gs}{\let\PY@bf=\textbf}
\@namedef{PY@tok@ges}{\let\PY@bf=\textbf\let\PY@it=\textit}
\@namedef{PY@tok@gp}{\let\PY@bf=\textbf\def\PY@tc##1{\textcolor[rgb]{0.00,0.00,0.50}{##1}}}
\@namedef{PY@tok@go}{\def\PY@tc##1{\textcolor[rgb]{0.44,0.44,0.44}{##1}}}
\@namedef{PY@tok@gt}{\def\PY@tc##1{\textcolor[rgb]{0.00,0.27,0.87}{##1}}}
\@namedef{PY@tok@err}{\def\PY@bc##1{{\setlength{\fboxsep}{\string -\fboxrule}\fcolorbox[rgb]{1.00,0.00,0.00}{1,1,1}{\strut ##1}}}}
\@namedef{PY@tok@kc}{\let\PY@bf=\textbf\def\PY@tc##1{\textcolor[rgb]{0.00,0.50,0.00}{##1}}}
\@namedef{PY@tok@kd}{\let\PY@bf=\textbf\def\PY@tc##1{\textcolor[rgb]{0.00,0.50,0.00}{##1}}}
\@namedef{PY@tok@kn}{\let\PY@bf=\textbf\def\PY@tc##1{\textcolor[rgb]{0.00,0.50,0.00}{##1}}}
\@namedef{PY@tok@kr}{\let\PY@bf=\textbf\def\PY@tc##1{\textcolor[rgb]{0.00,0.50,0.00}{##1}}}
\@namedef{PY@tok@bp}{\def\PY@tc##1{\textcolor[rgb]{0.00,0.50,0.00}{##1}}}
\@namedef{PY@tok@fm}{\def\PY@tc##1{\textcolor[rgb]{0.00,0.00,1.00}{##1}}}
\@namedef{PY@tok@vc}{\def\PY@tc##1{\textcolor[rgb]{0.10,0.09,0.49}{##1}}}
\@namedef{PY@tok@vg}{\def\PY@tc##1{\textcolor[rgb]{0.10,0.09,0.49}{##1}}}
\@namedef{PY@tok@vi}{\def\PY@tc##1{\textcolor[rgb]{0.10,0.09,0.49}{##1}}}
\@namedef{PY@tok@vm}{\def\PY@tc##1{\textcolor[rgb]{0.10,0.09,0.49}{##1}}}
\@namedef{PY@tok@sa}{\def\PY@tc##1{\textcolor[rgb]{0.73,0.13,0.13}{##1}}}
\@namedef{PY@tok@sb}{\def\PY@tc##1{\textcolor[rgb]{0.73,0.13,0.13}{##1}}}
\@namedef{PY@tok@sc}{\def\PY@tc##1{\textcolor[rgb]{0.73,0.13,0.13}{##1}}}
\@namedef{PY@tok@dl}{\def\PY@tc##1{\textcolor[rgb]{0.73,0.13,0.13}{##1}}}
\@namedef{PY@tok@s2}{\def\PY@tc##1{\textcolor[rgb]{0.73,0.13,0.13}{##1}}}
\@namedef{PY@tok@sh}{\def\PY@tc##1{\textcolor[rgb]{0.73,0.13,0.13}{##1}}}
\@namedef{PY@tok@s1}{\def\PY@tc##1{\textcolor[rgb]{0.73,0.13,0.13}{##1}}}
\@namedef{PY@tok@mb}{\def\PY@tc##1{\textcolor[rgb]{0.40,0.40,0.40}{##1}}}
\@namedef{PY@tok@mf}{\def\PY@tc##1{\textcolor[rgb]{0.40,0.40,0.40}{##1}}}
\@namedef{PY@tok@mh}{\def\PY@tc##1{\textcolor[rgb]{0.40,0.40,0.40}{##1}}}
\@namedef{PY@tok@mi}{\def\PY@tc##1{\textcolor[rgb]{0.40,0.40,0.40}{##1}}}
\@namedef{PY@tok@il}{\def\PY@tc##1{\textcolor[rgb]{0.40,0.40,0.40}{##1}}}
\@namedef{PY@tok@mo}{\def\PY@tc##1{\textcolor[rgb]{0.40,0.40,0.40}{##1}}}
\@namedef{PY@tok@ch}{\let\PY@it=\textit\def\PY@tc##1{\textcolor[rgb]{0.24,0.48,0.48}{##1}}}
\@namedef{PY@tok@cm}{\let\PY@it=\textit\def\PY@tc##1{\textcolor[rgb]{0.24,0.48,0.48}{##1}}}
\@namedef{PY@tok@cpf}{\let\PY@it=\textit\def\PY@tc##1{\textcolor[rgb]{0.24,0.48,0.48}{##1}}}
\@namedef{PY@tok@c1}{\let\PY@it=\textit\def\PY@tc##1{\textcolor[rgb]{0.24,0.48,0.48}{##1}}}
\@namedef{PY@tok@cs}{\let\PY@it=\textit\def\PY@tc##1{\textcolor[rgb]{0.24,0.48,0.48}{##1}}}

\def\PYZbs{\char`\\}
\def\PYZus{\char`\_}
\def\PYZob{\char`\{}
\def\PYZcb{\char`\}}
\def\PYZca{\char`\^}
\def\PYZam{\char`\&}
\def\PYZlt{\char`\<}
\def\PYZgt{\char`\>}
\def\PYZsh{\char`\#}
\def\PYZpc{\char`\%}
\def\PYZdl{\char`\$}
\def\PYZhy{\char`\-}
\def\PYZsq{\char`\'}
\def\PYZdq{\char`\"}
\def\PYZti{\char`\~}
% for compatibility with earlier versions
\def\PYZat{@}
\def\PYZlb{[}
\def\PYZrb{]}
\makeatother


% For linebreaks inside Verbatim environment from package fancyvrb.
\makeatletter
\newbox\Wrappedcontinuationbox
\newbox\Wrappedvisiblespacebox
\newcommand*\Wrappedvisiblespace {\textcolor{red}{\textvisiblespace}}
\newcommand*\Wrappedcontinuationsymbol {\textcolor{red}{\llap{\tiny$\m@th\hookrightarrow$}}}
\newcommand*\Wrappedcontinuationindent {3ex }
\newcommand*\Wrappedafterbreak {\kern\Wrappedcontinuationindent\copy\Wrappedcontinuationbox}
% Take advantage of the already applied Pygments mark-up to insert
% potential linebreaks for TeX processing.
%        {, <, #, %, $, ' and ": go to next line.
	%        _, }, ^, &, >, - and ~: stay at end of broken line.
% Use of \textquotesingle for straight quote.
\newcommand*\Wrappedbreaksatspecials {%
	\def\PYGZus{\discretionary{\char`\_}{\Wrappedafterbreak}{\char`\_}}%
	\def\PYGZob{\discretionary{}{\Wrappedafterbreak\char`\{}{\char`\{}}%
	\def\PYGZcb{\discretionary{\char`\}}{\Wrappedafterbreak}{\char`\}}}%
	\def\PYGZca{\discretionary{\char`\^}{\Wrappedafterbreak}{\char`\^}}%
	\def\PYGZam{\discretionary{\char`\&}{\Wrappedafterbreak}{\char`\&}}%
	\def\PYGZlt{\discretionary{}{\Wrappedafterbreak\char`\<}{\char`\<}}%
	\def\PYGZgt{\discretionary{\char`\>}{\Wrappedafterbreak}{\char`\>}}%
	\def\PYGZsh{\discretionary{}{\Wrappedafterbreak\char`\#}{\char`\#}}%
	\def\PYGZpc{\discretionary{}{\Wrappedafterbreak\char`\%}{\char`\%}}%
	\def\PYGZdl{\discretionary{}{\Wrappedafterbreak\char`\$}{\char`\$}}%
	\def\PYGZhy{\discretionary{\char`\-}{\Wrappedafterbreak}{\char`\-}}%
	\def\PYGZsq{\discretionary{}{\Wrappedafterbreak\textquotesingle}{\textquotesingle}}%
	\def\PYGZdq{\discretionary{}{\Wrappedafterbreak\char`\"}{\char`\"}}%
	\def\PYGZti{\discretionary{\char`\~}{\Wrappedafterbreak}{\char`\~}}%
}
% Some characters . , ; ? ! / are not pygmentized.
% This macro makes them "active" and they will insert potential linebreaks
\newcommand*\Wrappedbreaksatpunct {%
	\lccode`\~`\.\lowercase{\def~}{\discretionary{\hbox{\char`\.}}{\Wrappedafterbreak}{\hbox{\char`\.}}}%
	\lccode`\~`\,\lowercase{\def~}{\discretionary{\hbox{\char`\,}}{\Wrappedafterbreak}{\hbox{\char`\,}}}%
	\lccode`\~`\;\lowercase{\def~}{\discretionary{\hbox{\char`\;}}{\Wrappedafterbreak}{\hbox{\char`\;}}}%
	\lccode`\~`\:\lowercase{\def~}{\discretionary{\hbox{\char`\:}}{\Wrappedafterbreak}{\hbox{\char`\:}}}%
	\lccode`\~`\?\lowercase{\def~}{\discretionary{\hbox{\char`\?}}{\Wrappedafterbreak}{\hbox{\char`\?}}}%
	\lccode`\~`\!\lowercase{\def~}{\discretionary{\hbox{\char`\!}}{\Wrappedafterbreak}{\hbox{\char`\!}}}%
	\lccode`\~`\/\lowercase{\def~}{\discretionary{\hbox{\char`\/}}{\Wrappedafterbreak}{\hbox{\char`\/}}}%
	\catcode`\.\active
	\catcode`\,\active
	\catcode`\;\active
	\catcode`\:\active
	\catcode`\?\active
	\catcode`\!\active
	\catcode`\/\active
	\lccode`\~`\~
}
\makeatother

\let\OriginalVerbatim=\Verbatim
\makeatletter
\renewcommand{\Verbatim}[1][1]{%
	%\parskip\z@skip
	\sbox\Wrappedcontinuationbox {\Wrappedcontinuationsymbol}%
	\sbox\Wrappedvisiblespacebox {\FV@SetupFont\Wrappedvisiblespace}%
	\def\FancyVerbFormatLine ##1{\hsize\linewidth
		\vtop{\raggedright\hyphenpenalty\z@\exhyphenpenalty\z@
			\doublehyphendemerits\z@\finalhyphendemerits\z@
			\strut ##1\strut}%
	}%
	% If the linebreak is at a space, the latter will be displayed as visible
	% space at end of first line, and a continuation symbol starts next line.
	% Stretch/shrink are however usually zero for typewriter font.
	\def\FV@Space {%
		\nobreak\hskip\z@ plus\fontdimen3\font minus\fontdimen4\font
		\discretionary{\copy\Wrappedvisiblespacebox}{\Wrappedafterbreak}
		{\kern\fontdimen2\font}%
	}%
	
	% Allow breaks at special characters using \PYG... macros.
	\Wrappedbreaksatspecials
	% Breaks at punctuation characters . , ; ? ! and / need catcode=\active
	\OriginalVerbatim[#1,codes*=\Wrappedbreaksatpunct]%
}
\makeatother

% Exact colors from NB
\definecolor{incolor}{HTML}{303F9F}
\definecolor{outcolor}{HTML}{D84315}
\definecolor{cellborder}{HTML}{CFCFCF}
\definecolor{cellbackground}{HTML}{F7F7F7}

% prompt
\makeatletter
\newcommand{\boxspacing}{\kern\kvtcb@left@rule\kern\kvtcb@boxsep}
\makeatother
\newcommand{\prompt}[4]{
	{\ttfamily\llap{{\color{#2}[#3]:\hspace{3pt}#4}}\vspace{-\baselineskip}}
}



% Prevent overflowing lines due to hard-to-break entities
\sloppy
% Setup hyperref package
\hypersetup{
	breaklinks=true,  % so long urls are correctly broken across lines
	colorlinks=true,
	urlcolor=urlcolor,
	linkcolor=linkcolor,
	citecolor=citecolor,
}
% Slightly bigger margins than the latex defaults

\geometry{verbose,tmargin=1in,bmargin=1in,lmargin=1in,rmargin=1in}
    
    

\begin{document}
    
    \maketitle

	  \vspace*{4em}
	\begin{abstract}
		机器视觉第五次实验课程主要包含两大部分的练习，一是张正有标定法，二是SIFT算法。以十张照片为例，使用张正有标定法实现对图片的畸形矫正。对于SIFT算法部分，分别从两个不同的角度拍摄相同物品的照片，调用SIFT算法实现匹配。希望本篇文章是一篇内容翔实、方便随时翻阅复习学习OpenCV的学习笔记。
	\end{abstract}
	
	\begin{center}
		\vspace*{5em}{\large \textbf{Now the CODE is available at:}}\\\underline{https://github.com/GplasT0810/Machine-Vision-Experiment.git}\\
		\vspace*{6em}{\Large \textbf{长沙理工大学}}\\
		\vspace{0.7em}{\large \textbf{卓越工程师学院\ 绿色智慧交通 2401班}}\\
		\vspace{0.7em}\textbf{{\large 202404080608}}\\
		\vspace{0.7em}{\large \textbf{国嘉通}}
	\end{center}
	\newpage
	
	\tableofcontents
	\begin{center}
		\vspace{15em}\adjustimage{max size={0.9\linewidth}{0.9\paperheight}}{SchPic.png}
	\end{center}
	
	\vspace*{18em}--------------------------------------------\\
	Complated Time: Oct.$1^{st}$.2025\\ \\
	Guo Jiatong , Green \& Intelligent Transportation 2401\\
	Elite Engineering School , Changsha University of Science \& Technology.\\
	Changsha 410114 , Hunan , China
	
	\newpage
	
	
	\section{张正有标定法}
	张正友（Zhang Zhengyou）标定法，是 1998 年由微软研究院张正友提出的 单平面棋盘格相机标定法，也是目前计算机视觉领域使用最广泛的相机内参/外参标定技术。其可以仅用一张打印的棋盘格，在相机前任意摆几个姿势，就能高精度地求出相机内参、外参和畸变系数。但所打印的棋盘格必须全部暴露在相机视角下。
    \begin{tcolorbox}[breakable, size=fbox, boxrule=1pt, pad at break*=1mm,colback=cellbackground, colframe=cellborder]
%\prompt{In}{incolor}{7}{\boxspacing}
\begin{Verbatim}[commandchars=\\\{\}]
\PY{c+c1}{\PYZsh{}张正友标定}
\PY{k+kn}{import}\PY{+w}{ }\PY{n+nn}{cv2}
\PY{k+kn}{import}\PY{+w}{ }\PY{n+nn}{numpy}\PY{+w}{ }\PY{k}{as}\PY{+w}{ }\PY{n+nn}{np}
\PY{k+kn}{import}\PY{+w}{ }\PY{n+nn}{glob}
\PY{k+kn}{import}\PY{+w}{ }\PY{n+nn}{os}
\PY{k+kn}{import}\PY{+w}{ }\PY{n+nn}{matplotlib}\PY{n+nn}{.}\PY{n+nn}{pyplot}\PY{+w}{ }\PY{k}{as}\PY{+w}{ }\PY{n+nn}{plt}

\PY{c+c1}{\PYZsh{}对象点}
\PY{n}{chessboard\PYZus{}size} \PY{o}{=} \PY{p}{(}\PY{l+m+mi}{5}\PY{p}{,} \PY{l+m+mi}{8}\PY{p}{)}  \PY{c+c1}{\PYZsh{} 内角点数量（列，行）}
\PY{n}{square\PYZus{}size} \PY{o}{=} \PY{l+m+mf}{27.0}
\PY{c+c1}{\PYZsh{}实际标定使用的真实棋格的真实尺寸27mm}

\PY{n}{objp} \PY{o}{=} \PY{n}{np}\PY{o}{.}\PY{n}{zeros}\PY{p}{(}\PY{p}{(}\PY{n}{chessboard\PYZus{}size}\PY{p}{[}\PY{l+m+mi}{0}\PY{p}{]} \PY{o}{*} \PY{n}{chessboard\PYZus{}size}\PY{p}{[}\PY{l+m+mi}{1}\PY{p}{]}\PY{p}{,} \PY{l+m+mi}{3}\PY{p}{)}\PY{p}{,} \PY{n}{np}\PY{o}{.}\PY{n}{float32}\PY{p}{)}
\PY{n}{objp}\PY{p}{[}\PY{p}{:}\PY{p}{,} \PY{p}{:}\PY{l+m+mi}{2}\PY{p}{]} \PY{o}{=} \PY{n}{np}\PY{o}{.}\PY{n}{mgrid}\PY{p}{[}\PY{l+m+mi}{0}\PY{p}{:}\PY{n}{chessboard\PYZus{}size}\PY{p}{[}\PY{l+m+mi}{0}\PY{p}{]}\PY{p}{,} \PY{l+m+mi}{0}\PY{p}{:}\PY{n}{chessboard\PYZus{}size}\PY{p}{[}\PY{l+m+mi}{1}\PY{p}{]}\PY{p}{]}\PY{o}{.}\PY{n}{T}\PY{o}{.}\PY{n}{reshape}\PY{p}{(}\PY{o}{\PYZhy{}}\PY{l+m+mi}{1}\PY{p}{,} \PY{l+m+mi}{2}\PY{p}{)} \PY{o}{*} \PY{n}{square\PYZus{}size}
\PY{n}{objp}\PY{p}{[}\PY{p}{:}\PY{p}{,} \PY{p}{:}\PY{l+m+mi}{2}\PY{p}{]} \PY{o}{*}\PY{o}{=} \PY{n}{square\PYZus{}size}

\PY{c+c1}{\PYZsh{}存储对象点和图像点的数组}
\PY{n}{objpoints} \PY{o}{=} \PY{p}{[}\PY{p}{]}  \PY{c+c1}{\PYZsh{} 真实世界中的3D点}
\PY{n}{imgpoints} \PY{o}{=} \PY{p}{[}\PY{p}{]}  \PY{c+c1}{\PYZsh{} 图像中的2D点}

\PY{c+c1}{\PYZsh{}读取所有标定图像}
\PY{n}{folder} \PY{o}{=} \PY{l+s+sa}{r}\PY{l+s+s1}{\PYZsq{}}\PY{l+s+s1}{/Users/guo2006/myenv/Machine Vision Experiment/Machine Vision Experiment 5/biaoding}\PY{l+s+s1}{\PYZsq{}}
\PY{n}{images} \PY{o}{=} \PY{n}{glob}\PY{o}{.}\PY{n}{glob}\PY{p}{(}\PY{n}{os}\PY{o}{.}\PY{n}{path}\PY{o}{.}\PY{n}{join}\PY{p}{(}\PY{n}{folder}\PY{p}{,} \PY{l+s+s1}{\PYZsq{}}\PY{l+s+s1}{*.jpg}\PY{l+s+s1}{\PYZsq{}}\PY{p}{)}\PY{p}{)} \PYZbs{}
       \PY{o}{+} \PY{n}{glob}\PY{o}{.}\PY{n}{glob}\PY{p}{(}\PY{n}{os}\PY{o}{.}\PY{n}{path}\PY{o}{.}\PY{n}{join}\PY{p}{(}\PY{n}{folder}\PY{p}{,} \PY{l+s+s1}{\PYZsq{}}\PY{l+s+s1}{*.png}\PY{l+s+s1}{\PYZsq{}}\PY{p}{)}\PY{p}{)}  \PY{c+c1}{\PYZsh{} 两种后缀都读取}

\PY{k}{if} \PY{n+nb}{len}\PY{p}{(}\PY{n}{images}\PY{p}{)} \PY{o}{==} \PY{l+m+mi}{0}\PY{p}{:}
    \PY{n+nb}{print}\PY{p}{(}\PY{l+s+s2}{\PYZdq{}}\PY{l+s+s2}{未找到标定图像，请检查路径是否正确}\PY{l+s+s2}{\PYZdq{}}\PY{p}{)}
    \PY{n}{exit}\PY{p}{(}\PY{p}{)}

\PY{n+nb}{print}\PY{p}{(}\PY{l+s+sa}{f}\PY{l+s+s2}{\PYZdq{}}\PY{l+s+s2}{找到 }\PY{l+s+si}{\PYZob{}}\PY{n+nb}{len}\PY{p}{(}\PY{n}{images}\PY{p}{)}\PY{l+s+si}{\PYZcb{}}\PY{l+s+s2}{ 张标定图像}\PY{l+s+s2}{\PYZdq{}}\PY{p}{)}

\PY{c+c1}{\PYZsh{}遍历每张图像进行角点检测}
\PY{k}{for} \PY{n}{fname} \PY{o+ow}{in} \PY{n}{images}\PY{p}{:}
    \PY{n}{img} \PY{o}{=} \PY{n}{cv2}\PY{o}{.}\PY{n}{imread}\PY{p}{(}\PY{n}{fname}\PY{p}{)}
    \PY{k}{if} \PY{n}{img} \PY{o+ow}{is} \PY{k+kc}{None}\PY{p}{:}
        \PY{n+nb}{print}\PY{p}{(}\PY{l+s+sa}{f}\PY{l+s+s2}{\PYZdq{}}\PY{l+s+s2}{无法读取图像: }\PY{l+s+si}{\PYZob{}}\PY{n}{fname}\PY{l+s+si}{\PYZcb{}}\PY{l+s+s2}{\PYZdq{}}\PY{p}{)}
        \PY{k}{continue}
        
    \PY{n}{gray} \PY{o}{=} \PY{n}{cv2}\PY{o}{.}\PY{n}{cvtColor}\PY{p}{(}\PY{n}{img}\PY{p}{,} \PY{n}{cv2}\PY{o}{.}\PY{n}{COLOR\PYZus{}BGR2GRAY}\PY{p}{)}
    
    \PY{c+c1}{\PYZsh{} 查找棋格角点}
    \PY{n}{ret}\PY{p}{,} \PY{n}{corners} \PY{o}{=} \PY{n}{cv2}\PY{o}{.}\PY{n}{findChessboardCorners}\PY{p}{(}\PY{n}{gray}\PY{p}{,} \PY{n}{chessboard\PYZus{}size}\PY{p}{,} \PY{k+kc}{None}\PY{p}{)}
    
    \PY{k}{if} \PY{n}{ret}\PY{p}{:}
        \PY{n+nb}{print}\PY{p}{(}\PY{l+s+sa}{f}\PY{l+s+s2}{\PYZdq{}}\PY{l+s+s2}{在 }\PY{l+s+si}{\PYZob{}}\PY{n}{fname}\PY{l+s+si}{\PYZcb{}}\PY{l+s+s2}{ 中找到角点}\PY{l+s+s2}{\PYZdq{}}\PY{p}{)}
        \PY{n}{objpoints}\PY{o}{.}\PY{n}{append}\PY{p}{(}\PY{n}{objp}\PY{p}{)}
        
        \PY{c+c1}{\PYZsh{} 提高角点检测精度}
        \PY{n}{corners2} \PY{o}{=} \PY{n}{cv2}\PY{o}{.}\PY{n}{cornerSubPix}\PY{p}{(}\PY{n}{gray}\PY{p}{,} \PY{n}{corners}\PY{p}{,} \PY{p}{(}\PY{l+m+mi}{11}\PY{p}{,} \PY{l+m+mi}{11}\PY{p}{)}\PY{p}{,} \PY{p}{(}\PY{o}{\PYZhy{}}\PY{l+m+mi}{1}\PY{p}{,} \PY{o}{\PYZhy{}}\PY{l+m+mi}{1}\PY{p}{)}\PY{p}{,} 
                                    \PY{n}{criteria}\PY{o}{=}\PY{p}{(}\PY{n}{cv2}\PY{o}{.}\PY{n}{TERM\PYZus{}CRITERIA\PYZus{}EPS} \PY{o}{+} \PY{n}{cv2}\PY{o}{.}\PY{n}{TERM\PYZus{}CRITERIA\PYZus{}MAX\PYZus{}ITER}\PY{p}{,} \PY{l+m+mi}{30}\PY{p}{,} \PY{l+m+mf}{0.001}\PY{p}{)}\PY{p}{)}
        \PY{n}{imgpoints}\PY{o}{.}\PY{n}{append}\PY{p}{(}\PY{n}{corners2}\PY{p}{)}
        
        \PY{c+c1}{\PYZsh{} 可视化角点}
        \PY{c+c1}{\PYZsh{} img = cv2.drawChessboardCorners(img, chessboard\PYZus{}size, corners2, ret)}
        \PY{c+c1}{\PYZsh{} plt.imshow(img)}
        \PY{c+c1}{\PYZsh{} plt.title(\PYZsq{}Corners\PYZsq{})}
    \PY{k}{else}\PY{p}{:}
        \PY{n+nb}{print}\PY{p}{(}\PY{l+s+sa}{f}\PY{l+s+s2}{\PYZdq{}}\PY{l+s+s2}{在 }\PY{l+s+si}{\PYZob{}}\PY{n}{fname}\PY{l+s+si}{\PYZcb{}}\PY{l+s+s2}{ 中未找到角点}\PY{l+s+s2}{\PYZdq{}}\PY{p}{)}


\PY{c+c1}{\PYZsh{}检查是否找到足够的图像进行标定}
\PY{k}{if} \PY{n+nb}{len}\PY{p}{(}\PY{n}{objpoints}\PY{p}{)} \PY{o}{\PYZlt{}} \PY{l+m+mi}{3}\PY{p}{:}
    \PY{n+nb}{print}\PY{p}{(}\PY{l+s+s2}{\PYZdq{}}\PY{l+s+s2}{找到的可用图像太少，至少需要3张图像}\PY{l+s+s2}{\PYZdq{}}\PY{p}{)}
    \PY{n}{exit}\PY{p}{(}\PY{p}{)}

\PY{n+nb}{print}\PY{p}{(}\PY{l+s+sa}{f}\PY{l+s+s2}{\PYZdq{}}\PY{l+s+s2}{使用 }\PY{l+s+si}{\PYZob{}}\PY{n+nb}{len}\PY{p}{(}\PY{n}{objpoints}\PY{p}{)}\PY{l+s+si}{\PYZcb{}}\PY{l+s+s2}{ 张图像进行标定}\PY{l+s+s2}{\PYZdq{}}\PY{p}{)}

\PY{c+c1}{\PYZsh{}执行相机标定}
\PY{n+nb}{print}\PY{p}{(}\PY{l+s+s2}{\PYZdq{}}\PY{l+s+s2}{开始相机标定...}\PY{l+s+s2}{\PYZdq{}}\PY{p}{)}
\PY{n}{ret}\PY{p}{,} \PY{n}{mtx}\PY{p}{,} \PY{n}{dist}\PY{p}{,} \PY{n}{rvecs}\PY{p}{,} \PY{n}{tvecs} \PY{o}{=} \PY{n}{cv2}\PY{o}{.}\PY{n}{calibrateCamera}\PY{p}{(}
    \PY{n}{objpoints}\PY{p}{,} \PY{n}{imgpoints}\PY{p}{,} \PY{n}{gray}\PY{o}{.}\PY{n}{shape}\PY{p}{[}\PY{p}{:}\PY{p}{:}\PY{o}{\PYZhy{}}\PY{l+m+mi}{1}\PY{p}{]}\PY{p}{,} \PY{k+kc}{None}\PY{p}{,} \PY{k+kc}{None}
\PY{p}{)}

\PY{n+nb}{print}\PY{p}{(}\PY{l+s+s2}{\PYZdq{}}\PY{l+s+se}{\PYZbs{}n}\PY{l+s+s2}{标定结果:}\PY{l+s+s2}{\PYZdq{}}\PY{p}{)}
\PY{n+nb}{print}\PY{p}{(}\PY{l+s+sa}{f}\PY{l+s+s2}{\PYZdq{}}\PY{l+s+s2}{重投影误差: }\PY{l+s+si}{\PYZob{}}\PY{n}{ret}\PY{l+s+si}{\PYZcb{}}\PY{l+s+s2}{\PYZdq{}}\PY{p}{)}
\PY{n+nb}{print}\PY{p}{(}\PY{l+s+s2}{\PYZdq{}}\PY{l+s+s2}{相机内参矩阵:}\PY{l+s+s2}{\PYZdq{}}\PY{p}{)}
\PY{n+nb}{print}\PY{p}{(}\PY{n}{mtx}\PY{p}{)}
\PY{n+nb}{print}\PY{p}{(}\PY{l+s+s2}{\PYZdq{}}\PY{l+s+s2}{畸变系数:}\PY{l+s+s2}{\PYZdq{}}\PY{p}{)}
\PY{n+nb}{print}\PY{p}{(}\PY{n}{dist}\PY{p}{)}

\PY{c+c1}{\PYZsh{}标定结果验证 \PYZhy{} 计算重投影误差}
\PY{n}{mean\PYZus{}error} \PY{o}{=} \PY{l+m+mi}{0}
\PY{k}{for} \PY{n}{i} \PY{o+ow}{in} \PY{n+nb}{range}\PY{p}{(}\PY{n+nb}{len}\PY{p}{(}\PY{n}{objpoints}\PY{p}{)}\PY{p}{)}\PY{p}{:}
    \PY{n}{imgpoints2}\PY{p}{,} \PY{n}{\PYZus{}} \PY{o}{=} \PY{n}{cv2}\PY{o}{.}\PY{n}{projectPoints}\PY{p}{(}\PY{n}{objpoints}\PY{p}{[}\PY{n}{i}\PY{p}{]}\PY{p}{,} \PY{n}{rvecs}\PY{p}{[}\PY{n}{i}\PY{p}{]}\PY{p}{,} \PY{n}{tvecs}\PY{p}{[}\PY{n}{i}\PY{p}{]}\PY{p}{,} \PY{n}{mtx}\PY{p}{,} \PY{n}{dist}\PY{p}{)}
    \PY{n}{error} \PY{o}{=} \PY{n}{cv2}\PY{o}{.}\PY{n}{norm}\PY{p}{(}\PY{n}{imgpoints}\PY{p}{[}\PY{n}{i}\PY{p}{]}\PY{p}{,} \PY{n}{imgpoints2}\PY{p}{,} \PY{n}{cv2}\PY{o}{.}\PY{n}{NORM\PYZus{}L2}\PY{p}{)} \PY{o}{/} \PY{n+nb}{len}\PY{p}{(}\PY{n}{imgpoints2}\PY{p}{)}
    \PY{n}{mean\PYZus{}error} \PY{o}{+}\PY{o}{=} \PY{n}{error}

\PY{n+nb}{print}\PY{p}{(}\PY{l+s+sa}{f}\PY{l+s+s2}{\PYZdq{}}\PY{l+s+se}{\PYZbs{}n}\PY{l+s+s2}{平均重投影误差: }\PY{l+s+si}{\PYZob{}}\PY{n}{mean\PYZus{}error}\PY{+w}{ }\PY{o}{/}\PY{+w}{ }\PY{n+nb}{len}\PY{p}{(}\PY{n}{objpoints}\PY{p}{)}\PY{l+s+si}{\PYZcb{}}\PY{l+s+s2}{ 像素}\PY{l+s+s2}{\PYZdq{}}\PY{p}{)}

\PY{c+c1}{\PYZsh{}保存标定结果}
\PY{n}{np}\PY{o}{.}\PY{n}{savez}\PY{p}{(}\PY{l+s+s1}{\PYZsq{}}\PY{l+s+s1}{camera\PYZus{}calibration.npz}\PY{l+s+s1}{\PYZsq{}}\PY{p}{,} \PY{n}{mtx}\PY{o}{=}\PY{n}{mtx}\PY{p}{,} \PY{n}{dist}\PY{o}{=}\PY{n}{dist}\PY{p}{,} \PY{n}{rvecs}\PY{o}{=}\PY{n}{rvecs}\PY{p}{,} \PY{n}{tvecs}\PY{o}{=}\PY{n}{tvecs}\PY{p}{)}
\PY{n+nb}{print}\PY{p}{(}\PY{l+s+s2}{\PYZdq{}}\PY{l+s+se}{\PYZbs{}n}\PY{l+s+s2}{标定结果已保存到 camera\PYZus{}calibration.npz}\PY{l+s+s2}{\PYZdq{}}\PY{p}{)}

\PY{c+c1}{\PYZsh{}使用标定结果进行畸变校正}
\PY{k}{if} \PY{n+nb}{len}\PY{p}{(}\PY{n}{images}\PY{p}{)} \PY{o}{\PYZgt{}} \PY{l+m+mi}{0}\PY{p}{:}
    \PY{c+c1}{\PYZsh{} 使用第一张图像进行演示}
    \PY{n}{img} \PY{o}{=} \PY{n}{cv2}\PY{o}{.}\PY{n}{imread}\PY{p}{(}\PY{n}{images}\PY{p}{[}\PY{l+m+mi}{0}\PY{p}{]}\PY{p}{)}
    \PY{n}{h}\PY{p}{,} \PY{n}{w} \PY{o}{=} \PY{n}{img}\PY{o}{.}\PY{n}{shape}\PY{p}{[}\PY{p}{:}\PY{l+m+mi}{2}\PY{p}{]}
    
    \PY{c+c1}{\PYZsh{} 优化相机矩阵}
    \PY{n}{newcameramtx}\PY{p}{,} \PY{n}{roi} \PY{o}{=} \PY{n}{cv2}\PY{o}{.}\PY{n}{getOptimalNewCameraMatrix}\PY{p}{(}\PY{n}{mtx}\PY{p}{,} \PY{n}{dist}\PY{p}{,} \PY{p}{(}\PY{n}{w}\PY{p}{,} \PY{n}{h}\PY{p}{)}\PY{p}{,} \PY{l+m+mi}{1}\PY{p}{,} \PY{p}{(}\PY{n}{w}\PY{p}{,} \PY{n}{h}\PY{p}{)}\PY{p}{)}
    
    \PY{c+c1}{\PYZsh{} 校正图像}
    \PY{n}{dst} \PY{o}{=} \PY{n}{cv2}\PY{o}{.}\PY{n}{undistort}\PY{p}{(}\PY{n}{img}\PY{p}{,} \PY{n}{mtx}\PY{p}{,} \PY{n}{dist}\PY{p}{,} \PY{k+kc}{None}\PY{p}{,} \PY{n}{newcameramtx}\PY{p}{)}
    
    \PY{c+c1}{\PYZsh{} 裁剪图像}
    \PY{n}{x}\PY{p}{,} \PY{n}{y}\PY{p}{,} \PY{n}{w}\PY{p}{,} \PY{n}{h} \PY{o}{=} \PY{n}{roi}
    \PY{n}{dst} \PY{o}{=} \PY{n}{dst}\PY{p}{[}\PY{n}{y}\PY{p}{:}\PY{n}{y}\PY{o}{+}\PY{n}{h}\PY{p}{,} \PY{n}{x}\PY{p}{:}\PY{n}{x}\PY{o}{+}\PY{n}{w}\PY{p}{]}
    
    \PY{c+c1}{\PYZsh{} 显示结果}
    \PY{n}{fig}\PY{p}{,} \PY{n}{ax} \PY{o}{=} \PY{n}{plt}\PY{o}{.}\PY{n}{subplots}\PY{p}{(}\PY{l+m+mi}{1}\PY{p}{,} \PY{l+m+mi}{2}\PY{p}{,} \PY{n}{figsize}\PY{o}{=}\PY{p}{(}\PY{l+m+mi}{10}\PY{p}{,}\PY{l+m+mi}{8}\PY{p}{)}\PY{p}{)}
    \PY{n}{ax}\PY{p}{[}\PY{l+m+mi}{0}\PY{p}{]}\PY{o}{.}\PY{n}{imshow}\PY{p}{(}\PY{n}{img}\PY{p}{)}\PY{p}{,} \PY{n}{ax}\PY{p}{[}\PY{l+m+mi}{0}\PY{p}{]}\PY{o}{.}\PY{n}{set\PYZus{}title}\PY{p}{(}\PY{l+s+s1}{\PYZsq{}}\PY{l+s+s1}{Original Image}\PY{l+s+s1}{\PYZsq{}}\PY{p}{)}\PY{p}{,} \PY{n}{ax}\PY{p}{[}\PY{l+m+mi}{0}\PY{p}{]}\PY{o}{.}\PY{n}{axis}\PY{p}{(}\PY{l+s+s1}{\PYZsq{}}\PY{l+s+s1}{off}\PY{l+s+s1}{\PYZsq{}}\PY{p}{)}
    \PY{n}{ax}\PY{p}{[}\PY{l+m+mi}{1}\PY{p}{]}\PY{o}{.}\PY{n}{imshow}\PY{p}{(}\PY{n}{dst}\PY{p}{)}\PY{p}{,} \PY{n}{ax}\PY{p}{[}\PY{l+m+mi}{1}\PY{p}{]}\PY{o}{.}\PY{n}{set\PYZus{}title}\PY{p}{(}\PY{l+s+s1}{\PYZsq{}}\PY{l+s+s1}{Undistorted Image}\PY{l+s+s1}{\PYZsq{}}\PY{p}{)}\PY{p}{,} \PY{n}{ax}\PY{p}{[}\PY{l+m+mi}{1}\PY{p}{]}\PY{o}{.}\PY{n}{axis}\PY{p}{(}\PY{l+s+s1}{\PYZsq{}}\PY{l+s+s1}{off}\PY{l+s+s1}{\PYZsq{}}\PY{p}{)}

    \PY{c+c1}{\PYZsh{} \PYZhy{}\PYZhy{}\PYZhy{}\PYZhy{}\PYZhy{}\PYZhy{}\PYZhy{}\PYZhy{}\PYZhy{}\PYZhy{} 带角点的畸变校正 \PYZhy{}\PYZhy{}\PYZhy{}\PYZhy{}\PYZhy{}\PYZhy{}\PYZhy{}\PYZhy{}\PYZhy{}\PYZhy{}}
\PY{k}{if} \PY{n+nb}{len}\PY{p}{(}\PY{n}{images}\PY{p}{)} \PY{o}{\PYZgt{}} \PY{l+m+mi}{0}\PY{p}{:}
    \PY{n}{img}  \PY{o}{=} \PY{n}{cv2}\PY{o}{.}\PY{n}{imread}\PY{p}{(}\PY{n}{images}\PY{p}{[}\PY{l+m+mi}{0}\PY{p}{]}\PY{p}{)}
    \PY{n}{gray} \PY{o}{=} \PY{n}{cv2}\PY{o}{.}\PY{n}{cvtColor}\PY{p}{(}\PY{n}{img}\PY{p}{,} \PY{n}{cv2}\PY{o}{.}\PY{n}{COLOR\PYZus{}BGR2GRAY}\PY{p}{)}
    \PY{n}{h}\PY{p}{,} \PY{n}{w} \PY{o}{=} \PY{n}{img}\PY{o}{.}\PY{n}{shape}\PY{p}{[}\PY{p}{:}\PY{l+m+mi}{2}\PY{p}{]}

    \PY{c+c1}{\PYZsh{} （1）先在校正前图上重新检角点并画上去}
    \PY{n}{ret}\PY{p}{,} \PY{n}{corners} \PY{o}{=} \PY{n}{cv2}\PY{o}{.}\PY{n}{findChessboardCorners}\PY{p}{(}\PY{n}{gray}\PY{p}{,} \PY{n}{chessboard\PYZus{}size}\PY{p}{,}
                                             \PY{n}{cv2}\PY{o}{.}\PY{n}{CALIB\PYZus{}CB\PYZus{}ADAPTIVE\PYZus{}THRESH} \PY{o}{+}
                                             \PY{n}{cv2}\PY{o}{.}\PY{n}{CALIB\PYZus{}CB\PYZus{}NORMALIZE\PYZus{}IMAGE}\PY{p}{)}
    \PY{k}{if} \PY{n}{ret}\PY{p}{:}                       \PY{c+c1}{\PYZsh{} 如果找到就画}
        \PY{n}{corners} \PY{o}{=} \PY{n}{cv2}\PY{o}{.}\PY{n}{cornerSubPix}\PY{p}{(}\PY{n}{gray}\PY{p}{,} \PY{n}{corners}\PY{p}{,} \PY{p}{(}\PY{l+m+mi}{11}\PY{p}{,} \PY{l+m+mi}{11}\PY{p}{)}\PY{p}{,} \PY{p}{(}\PY{o}{\PYZhy{}}\PY{l+m+mi}{1}\PY{p}{,} \PY{o}{\PYZhy{}}\PY{l+m+mi}{1}\PY{p}{)}\PY{p}{,}
                                   \PY{p}{(}\PY{n}{cv2}\PY{o}{.}\PY{n}{TERM\PYZus{}CRITERIA\PYZus{}EPS} \PY{o}{+} \PY{n}{cv2}\PY{o}{.}\PY{n}{TERM\PYZus{}CRITERIA\PYZus{}MAX\PYZus{}ITER}\PY{p}{,} \PY{l+m+mi}{30}\PY{p}{,} \PY{l+m+mf}{0.001}\PY{p}{)}\PY{p}{)}
        \PY{n}{cv2}\PY{o}{.}\PY{n}{drawChessboardCorners}\PY{p}{(}\PY{n}{img}\PY{p}{,} \PY{n}{chessboard\PYZus{}size}\PY{p}{,} \PY{n}{corners}\PY{p}{,} \PY{k+kc}{True}\PY{p}{)}

    \PY{c+c1}{\PYZsh{} （2）再去畸变}
    \PY{n}{newcameramtx}\PY{p}{,} \PY{n}{roi} \PY{o}{=} \PY{n}{cv2}\PY{o}{.}\PY{n}{getOptimalNewCameraMatrix}\PY{p}{(}\PY{n}{mtx}\PY{p}{,} \PY{n}{dist}\PY{p}{,} \PY{p}{(}\PY{n}{w}\PY{p}{,} \PY{n}{h}\PY{p}{)}\PY{p}{,} \PY{l+m+mi}{1}\PY{p}{,} \PY{p}{(}\PY{n}{w}\PY{p}{,} \PY{n}{h}\PY{p}{)}\PY{p}{)}
    \PY{n}{dst} \PY{o}{=} \PY{n}{cv2}\PY{o}{.}\PY{n}{undistort}\PY{p}{(}\PY{n}{img}\PY{p}{,} \PY{n}{mtx}\PY{p}{,} \PY{n}{dist}\PY{p}{,} \PY{k+kc}{None}\PY{p}{,} \PY{n}{newcameramtx}\PY{p}{)}

    \PY{c+c1}{\PYZsh{} （3）裁剪黑边（可选）}
    \PY{n}{x}\PY{p}{,} \PY{n}{y}\PY{p}{,} \PY{n}{w}\PY{p}{,} \PY{n}{h} \PY{o}{=} \PY{n}{roi}
    \PY{n}{dst} \PY{o}{=} \PY{n}{dst}\PY{p}{[}\PY{n}{y}\PY{p}{:}\PY{n}{y}\PY{o}{+}\PY{n}{h}\PY{p}{,} \PY{n}{x}\PY{p}{:}\PY{n}{x}\PY{o}{+}\PY{n}{w}\PY{p}{]}

    \PY{c+c1}{\PYZsh{} （4）显示 \PYZam{} 保存}
    \PY{n}{fig}\PY{p}{,} \PY{n}{ax} \PY{o}{=} \PY{n}{plt}\PY{o}{.}\PY{n}{subplots}\PY{p}{(}\PY{l+m+mi}{1}\PY{p}{,} \PY{l+m+mi}{2}\PY{p}{,} \PY{n}{figsize}\PY{o}{=}\PY{p}{(}\PY{l+m+mi}{10}\PY{p}{,}\PY{l+m+mi}{8}\PY{p}{)}\PY{p}{)}
    \PY{n}{ax}\PY{p}{[}\PY{l+m+mi}{0}\PY{p}{]}\PY{o}{.}\PY{n}{imshow}\PY{p}{(}\PY{n}{img}\PY{p}{)}\PY{p}{,} \PY{n}{ax}\PY{p}{[}\PY{l+m+mi}{0}\PY{p}{]}\PY{o}{.}\PY{n}{set\PYZus{}title}\PY{p}{(}\PY{l+s+s1}{\PYZsq{}}\PY{l+s+s1}{Original Image}\PY{l+s+s1}{\PYZsq{}}\PY{p}{)}\PY{p}{,} \PY{n}{ax}\PY{p}{[}\PY{l+m+mi}{0}\PY{p}{]}\PY{o}{.}\PY{n}{axis}\PY{p}{(}\PY{l+s+s1}{\PYZsq{}}\PY{l+s+s1}{off}\PY{l+s+s1}{\PYZsq{}}\PY{p}{)}
    \PY{n}{ax}\PY{p}{[}\PY{l+m+mi}{1}\PY{p}{]}\PY{o}{.}\PY{n}{imshow}\PY{p}{(}\PY{n}{dst}\PY{p}{)}\PY{p}{,} \PY{n}{ax}\PY{p}{[}\PY{l+m+mi}{1}\PY{p}{]}\PY{o}{.}\PY{n}{set\PYZus{}title}\PY{p}{(}\PY{l+s+s1}{\PYZsq{}}\PY{l+s+s1}{Undistorted Image}\PY{l+s+s1}{\PYZsq{}}\PY{p}{)}\PY{p}{,} \PY{n}{ax}\PY{p}{[}\PY{l+m+mi}{1}\PY{p}{]}\PY{o}{.}\PY{n}{axis}\PY{p}{(}\PY{l+s+s1}{\PYZsq{}}\PY{l+s+s1}{off}\PY{l+s+s1}{\PYZsq{}}\PY{p}{)}

    \PY{c+c1}{\PYZsh{} 生成未校正 \PYZam{} 校正两张图（同尺寸，不裁剪）}
    \PY{n}{img\PYZus{}raw}   \PY{o}{=} \PY{n}{cv2}\PY{o}{.}\PY{n}{imread}\PY{p}{(}\PY{n}{images}\PY{p}{[}\PY{l+m+mi}{0}\PY{p}{]}\PY{p}{)}
    \PY{n}{h}\PY{p}{,} \PY{n}{w} \PY{o}{=} \PY{n}{img\PYZus{}raw}\PY{o}{.}\PY{n}{shape}\PY{p}{[}\PY{p}{:}\PY{l+m+mi}{2}\PY{p}{]}
    \PY{n}{map1}\PY{p}{,} \PY{n}{map2} \PY{o}{=} \PY{n}{cv2}\PY{o}{.}\PY{n}{initUndistortRectifyMap}\PY{p}{(}\PY{n}{mtx}\PY{p}{,} \PY{n}{dist}\PY{p}{,} \PY{k+kc}{None}\PY{p}{,} \PY{n}{mtx}\PY{p}{,} \PY{p}{(}\PY{n}{w}\PY{p}{,} \PY{n}{h}\PY{p}{)}\PY{p}{,} \PY{n}{cv2}\PY{o}{.}\PY{n}{CV\PYZus{}16SC2}\PY{p}{)}
    \PY{n}{img\PYZus{}undist} \PY{o}{=} \PY{n}{cv2}\PY{o}{.}\PY{n}{remap}\PY{p}{(}\PY{n}{img\PYZus{}raw}\PY{p}{,} \PY{n}{map1}\PY{p}{,} \PY{n}{map2}\PY{p}{,} \PY{n}{cv2}\PY{o}{.}\PY{n}{INTER\PYZus{}LINEAR}\PY{p}{)}

    \PY{c+c1}{\PYZsh{} 计算差异并增强对比}
    \PY{n}{diff} \PY{o}{=} \PY{n}{cv2}\PY{o}{.}\PY{n}{absdiff}\PY{p}{(}\PY{n}{img\PYZus{}raw}\PY{p}{,} \PY{n}{img\PYZus{}undist}\PY{p}{)}
    \PY{n}{diff} \PY{o}{=} \PY{n}{cv2}\PY{o}{.}\PY{n}{convertScaleAbs}\PY{p}{(}\PY{n}{diff}\PY{p}{,} \PY{n}{alpha}\PY{o}{=}\PY{l+m+mf}{1.5}\PY{p}{,} \PY{n}{beta}\PY{o}{=}\PY{l+m+mi}{0}\PY{p}{)}   \PY{c+c1}{\PYZsh{} 放大 5 倍}
    \PY{n}{plt}\PY{o}{.}\PY{n}{imshow}\PY{p}{(}\PY{n}{diff}\PY{p}{)}
    \PY{n}{plt}\PY{o}{.}\PY{n}{title}\PY{p}{(}\PY{l+s+s1}{\PYZsq{}}\PY{l+s+s1}{diff img}\PY{l+s+s1}{\PYZsq{}}\PY{p}{)}

    \PY{n+nb}{print}\PY{p}{(}\PY{l+s+s2}{\PYZdq{}}\PY{l+s+se}{\PYZbs{}n}\PY{l+s+s2}{标定完成！}\PY{l+s+s2}{\PYZdq{}}\PY{p}{)}
\end{Verbatim}
\end{tcolorbox}

    \begin{Verbatim}[commandchars=\\\{\}]
找到 20 张标定图像
在 /Users/guo2006/myenv/Machine Vision Experiment/Machine Vision Experiment
5/biaoding/8.jpg 中找到角点
在 /Users/guo2006/myenv/Machine Vision Experiment/Machine Vision Experiment
5/biaoding/9.jpg 中找到角点
在 /Users/guo2006/myenv/Machine Vision Experiment/Machine Vision Experiment
5/biaoding/14.jpg 中找到角点
在 /Users/guo2006/myenv/Machine Vision Experiment/Machine Vision Experiment
5/biaoding/15.jpg 中找到角点
在 /Users/guo2006/myenv/Machine Vision Experiment/Machine Vision Experiment
5/biaoding/17.jpg 中找到角点
在 /Users/guo2006/myenv/Machine Vision Experiment/Machine Vision Experiment
5/biaoding/16.jpg 中找到角点
在 /Users/guo2006/myenv/Machine Vision Experiment/Machine Vision Experiment
5/biaoding/12.jpg 中找到角点
在 /Users/guo2006/myenv/Machine Vision Experiment/Machine Vision Experiment
5/biaoding/13.jpg 中找到角点
在 /Users/guo2006/myenv/Machine Vision Experiment/Machine Vision Experiment
5/biaoding/11.jpg 中找到角点
在 /Users/guo2006/myenv/Machine Vision Experiment/Machine Vision Experiment
5/biaoding/10.jpg 中找到角点
在 /Users/guo2006/myenv/Machine Vision Experiment/Machine Vision Experiment
5/biaoding/20.jpg 中找到角点
在 /Users/guo2006/myenv/Machine Vision Experiment/Machine Vision Experiment
5/biaoding/18.jpg 中找到角点
在 /Users/guo2006/myenv/Machine Vision Experiment/Machine Vision Experiment
5/biaoding/19.jpg 中找到角点
在 /Users/guo2006/myenv/Machine Vision Experiment/Machine Vision Experiment
5/biaoding/4.jpg 中找到角点
在 /Users/guo2006/myenv/Machine Vision Experiment/Machine Vision Experiment
5/biaoding/5.jpg 中找到角点
在 /Users/guo2006/myenv/Machine Vision Experiment/Machine Vision Experiment
5/biaoding/7.jpg 中找到角点
在 /Users/guo2006/myenv/Machine Vision Experiment/Machine Vision Experiment
5/biaoding/6.jpg 中找到角点
在 /Users/guo2006/myenv/Machine Vision Experiment/Machine Vision Experiment
5/biaoding/2.jpg 中找到角点
在 /Users/guo2006/myenv/Machine Vision Experiment/Machine Vision Experiment
5/biaoding/3.jpg 中找到角点
在 /Users/guo2006/myenv/Machine Vision Experiment/Machine Vision Experiment
5/biaoding/1.jpg 中找到角点
使用 20 张图像进行标定
开始相机标定{\ldots}

标定结果:
重投影误差: 8.203850256183376
相机内参矩阵:
[[6.24862606e+03 0.00000000e+00 4.93473537e+02]
 [0.00000000e+00 1.82304138e+04 9.05418068e+02]
 [0.00000000e+00 0.00000000e+00 1.00000000e+00]]
畸变系数:
[[-4.27055210e-01  6.06555828e+00 -6.31486059e-03 -4.73005413e-02
  -2.35297951e+01]]

平均重投影误差: 1.0953467595808408 像素

标定结果已保存到 camera\_calibration.npz

标定完成！
    \end{Verbatim}

    \begin{center}
    \adjustimage{max size={0.9\linewidth}{0.9\paperheight}}{output_0_1.png}
    \end{center}
    
    \begin{center}
    \adjustimage{max size={0.9\linewidth}{0.9\paperheight}}{output_0_2.png}
    \end{center}
在“diff img”图中，白色部分代表畸形矫正后与畸形矫正前的变化。
    
    \section{SIFT算法}
    尺度不变特征转换(Scale-invariant feature transform 或 SIFT)是一种机器视觉的算法用来侦测与描述影像中的局部性特征，它在空间尺度中寻找极值点，并提取出其位置、尺度、旋转不变数，此算法由 David Lowe 在1999年所发表，2004年完善总结。 
    \begin{tcolorbox}[breakable, size=fbox, boxrule=1pt, pad at break*=1mm,colback=cellbackground, colframe=cellborder]
%\prompt{In}{incolor}{12}{\boxspacing}
\begin{Verbatim}[commandchars=\\\{\}]
\PY{k+kn}{import}\PY{+w}{ }\PY{n+nn}{cv2}
\PY{k+kn}{import}\PY{+w}{ }\PY{n+nn}{numpy}\PY{+w}{ }\PY{k}{as}\PY{+w}{ }\PY{n+nn}{np}
\PY{k+kn}{import}\PY{+w}{ }\PY{n+nn}{matplotlib}\PY{n+nn}{.}\PY{n+nn}{pyplot}\PY{+w}{ }\PY{k}{as}\PY{+w}{ }\PY{n+nn}{plt}

\PY{c+c1}{\PYZsh{} 读图}
\PY{n}{img1} \PY{o}{=} \PY{n}{cv2}\PY{o}{.}\PY{n}{imread}\PY{p}{(}\PY{l+s+s1}{\PYZsq{}}\PY{l+s+s1}{Ang1.jpg}\PY{l+s+s1}{\PYZsq{}}\PY{p}{,} \PY{n}{cv2}\PY{o}{.}\PY{n}{IMREAD\PYZus{}GRAYSCALE}\PY{p}{)}  \PY{c+c1}{\PYZsh{} 查询图}
\PY{n}{img2} \PY{o}{=} \PY{n}{cv2}\PY{o}{.}\PY{n}{imread}\PY{p}{(}\PY{l+s+s1}{\PYZsq{}}\PY{l+s+s1}{Ang2.jpg}\PY{l+s+s1}{\PYZsq{}}\PY{p}{,} \PY{n}{cv2}\PY{o}{.}\PY{n}{IMREAD\PYZus{}GRAYSCALE}\PY{p}{)}  \PY{c+c1}{\PYZsh{} 训练图}
\PY{k}{assert} \PY{n}{img1} \PY{o+ow}{is} \PY{o+ow}{not} \PY{k+kc}{None} \PY{o+ow}{and} \PY{n}{img2} \PY{o+ow}{is} \PY{o+ow}{not} \PY{k+kc}{None}\PY{p}{,} \PY{l+s+s2}{\PYZdq{}}\PY{l+s+s2}{无法读取图像，请检查路径！}\PY{l+s+s2}{\PYZdq{}}

\PY{c+c1}{\PYZsh{} 创建 SIFT 检测器}
\PY{n}{sift} \PY{o}{=} \PY{n}{cv2}\PY{o}{.}\PY{n}{SIFT\PYZus{}create}\PY{p}{(}\PY{p}{)}

\PY{c+c1}{\PYZsh{} 检测关键点和描述子}
\PY{n}{kp1}\PY{p}{,} \PY{n}{des1} \PY{o}{=} \PY{n}{sift}\PY{o}{.}\PY{n}{detectAndCompute}\PY{p}{(}\PY{n}{img1}\PY{p}{,} \PY{k+kc}{None}\PY{p}{)}
\PY{n}{kp2}\PY{p}{,} \PY{n}{des2} \PY{o}{=} \PY{n}{sift}\PY{o}{.}\PY{n}{detectAndCompute}\PY{p}{(}\PY{n}{img2}\PY{p}{,} \PY{k+kc}{None}\PY{p}{)}
\PY{n+nb}{print}\PY{p}{(}\PY{l+s+sa}{f}\PY{l+s+s2}{\PYZdq{}}\PY{l+s+s2}{图1关键点: }\PY{l+s+si}{\PYZob{}}\PY{n+nb}{len}\PY{p}{(}\PY{n}{kp1}\PY{p}{)}\PY{l+s+si}{\PYZcb{}}\PY{l+s+s2}{ 个，图2关键点: }\PY{l+s+si}{\PYZob{}}\PY{n+nb}{len}\PY{p}{(}\PY{n}{kp2}\PY{p}{)}\PY{l+s+si}{\PYZcb{}}\PY{l+s+s2}{ 个}\PY{l+s+s2}{\PYZdq{}}\PY{p}{)}

\PY{c+c1}{\PYZsh{} 暴力匹配 + 比值测试}
\PY{n}{bf} \PY{o}{=} \PY{n}{cv2}\PY{o}{.}\PY{n}{BFMatcher}\PY{p}{(}\PY{n}{cv2}\PY{o}{.}\PY{n}{NORM\PYZus{}L2}\PY{p}{,} \PY{n}{crossCheck}\PY{o}{=}\PY{k+kc}{False}\PY{p}{)}
\PY{n}{matches} \PY{o}{=} \PY{n}{bf}\PY{o}{.}\PY{n}{knnMatch}\PY{p}{(}\PY{n}{des1}\PY{p}{,} \PY{n}{des2}\PY{p}{,} \PY{n}{k}\PY{o}{=}\PY{l+m+mi}{2}\PY{p}{)}

\PY{n}{good} \PY{o}{=} \PY{p}{[}\PY{p}{]}
\PY{k}{for} \PY{n}{m}\PY{p}{,} \PY{n}{n} \PY{o+ow}{in} \PY{n}{matches}\PY{p}{:}
    \PY{c+c1}{\PYZsh{} Lowe 比值测试：最近距离 \PYZlt{} 0.7*次近距离}
    \PY{k}{if} \PY{n}{m}\PY{o}{.}\PY{n}{distance} \PY{o}{\PYZlt{}} \PY{l+m+mf}{0.7} \PY{o}{*} \PY{n}{n}\PY{o}{.}\PY{n}{distance}\PY{p}{:}
        \PY{n}{good}\PY{o}{.}\PY{n}{append}\PY{p}{(}\PY{n}{m}\PY{p}{)}
\PY{n+nb}{print}\PY{p}{(}\PY{l+s+s2}{\PYZdq{}}\PY{l+s+s2}{良好匹配对数：}\PY{l+s+s2}{\PYZdq{}}\PY{p}{,} \PY{n+nb}{len}\PY{p}{(}\PY{n}{good}\PY{p}{)}\PY{p}{)}

\PY{c+c1}{\PYZsh{} 可视化}
\PY{n}{img\PYZus{}match} \PY{o}{=} \PY{n}{cv2}\PY{o}{.}\PY{n}{drawMatches}\PY{p}{(}
    \PY{n}{img1}\PY{p}{,} \PY{n}{kp1}\PY{p}{,} \PY{n}{img2}\PY{p}{,} \PY{n}{kp2}\PY{p}{,} \PY{n}{good}\PY{p}{,} \PY{k+kc}{None}\PY{p}{,}
    \PY{n}{flags}\PY{o}{=}\PY{n}{cv2}\PY{o}{.}\PY{n}{DrawMatchesFlags\PYZus{}NOT\PYZus{}DRAW\PYZus{}SINGLE\PYZus{}POINTS}\PY{p}{)}

\PY{n}{plt}\PY{o}{.}\PY{n}{figure}\PY{p}{(}\PY{n}{figsize}\PY{o}{=}\PY{p}{(}\PY{l+m+mi}{12}\PY{p}{,} \PY{l+m+mi}{6}\PY{p}{)}\PY{p}{)}
\PY{n}{plt}\PY{o}{.}\PY{n}{imshow}\PY{p}{(}\PY{n}{img\PYZus{}match}\PY{p}{[}\PY{p}{:}\PY{p}{,} \PY{p}{:}\PY{p}{,} \PY{p}{:}\PY{p}{:}\PY{o}{\PYZhy{}}\PY{l+m+mi}{1}\PY{p}{]}\PY{p}{)}  \PY{c+c1}{\PYZsh{} BGR→RGB}
\PY{n}{plt}\PY{o}{.}\PY{n}{title}\PY{p}{(}\PY{l+s+sa}{f}\PY{l+s+s2}{\PYZdq{}}\PY{l+s+s2}{SIFT matches (After Ratio Test): }\PY{l+s+si}{\PYZob{}}\PY{n+nb}{len}\PY{p}{(}\PY{n}{good}\PY{p}{)}\PY{l+s+si}{\PYZcb{}}\PY{l+s+s2}{ Points}\PY{l+s+s2}{\PYZdq{}}\PY{p}{)}
\PY{n}{plt}\PY{o}{.}\PY{n}{axis}\PY{p}{(}\PY{l+s+s1}{\PYZsq{}}\PY{l+s+s1}{off}\PY{l+s+s1}{\PYZsq{}}\PY{p}{)}
\PY{n}{plt}\PY{o}{.}\PY{n}{tight\PYZus{}layout}\PY{p}{(}\PY{p}{)}
\PY{n}{plt}\PY{o}{.}\PY{n}{show}\PY{p}{(}\PY{p}{)}
\end{Verbatim}
\end{tcolorbox}

    \begin{Verbatim}[commandchars=\\\{\}]
图1关键点: 5376 个，图2关键点: 7051 个
良好匹配对数： 101
    \end{Verbatim}

    \begin{center}
    \adjustimage{max size={0.9\linewidth}{0.9\paperheight}}{output_1_1.png}
    \end{center}


    % Add a bibliography block to the postdoc
    
    
    
\end{document}
